% Step 4 proposals: mitigating correction avoidance.
% Compiles with pdfLaTeX on Overleaf (standard TeX Live packages only).
\documentclass[11pt]{article}

\usepackage[T1]{fontenc}
\usepackage[utf8]{inputenc}
\usepackage{lmodern}
\usepackage[margin=0.95in]{geometry}
\usepackage{microtype}
\usepackage{amsmath,amssymb}
\usepackage{booktabs,tabularx,array}
\usepackage[shortlabels]{enumitem}
\usepackage[dvipsnames]{xcolor}
\usepackage{tcolorbox}
\tcbuselibrary{breakable,skins}
\usepackage{titlesec}
\usepackage[colorlinks=true,linkcolor=accent,citecolor=accent,urlcolor=accent]{hyperref}

% ---------- style ----------
\definecolor{accent}{HTML}{1F4E79}
\definecolor{soft}{HTML}{EEF3F8}
\definecolor{muted}{HTML}{5A6470}

\titleformat{\section}{\large\bfseries\color{accent}}{\thesection}{0.6em}{}
\titlespacing*{\section}{0pt}{1.1em}{0.45em}

\setlength{\parindent}{0pt}
\setlength{\parskip}{0.45em}
\setlist{nosep,leftmargin=1.3em}
\newcolumntype{L}{>{\raggedright\arraybackslash}X}

\newtcolorbox{keybox}{enhanced, colback=soft, colframe=soft, arc=2pt,
  left=8pt, right=8pt, top=6pt, bottom=6pt}
\newtcolorbox{proposal}[1]{enhanced, breakable, colback=white, colframe=accent,
  boxrule=0.6pt, arc=2pt, left=8pt, right=8pt, top=5pt, bottom=6pt,
  fonttitle=\bfseries, coltitle=white, title={#1}}

\newcommand{\fld}[1]{\par\vspace{3pt}{\color{accent}\textbf{#1.}}\enspace}
\newcommand{\src}[1]{{\footnotesize\textsf{\color{muted}[#1]}}}
\newcommand{\lever}[1]{\colorbox{soft}{\footnotesize\textsf{lever: #1}}}
\newcommand{\code}[1]{\texttt{\small #1}}

% =====================================================================
\begin{document}

\begin{center}
  {\LARGE\bfseries\color{accent} Step 4: Stopping the Agent from Avoiding Correction}\\[5pt]
  {\large Proposals and experiment sketches}\\[6pt]
  {\color{muted} Cody Qiu \quad$\cdot$\quad Draft, 25 October 2026}
\end{center}

\begin{keybox}
\textbf{Goal.} During training the agent exploits a reward bug outside the
audit window and complies inside it, so it is never corrected. Step 4 looks for
ways to stop this. Any end state other than hiding is progress:
\par\smallskip
{\small
\begin{tabularx}{\linewidth}{@{}lL@{}}
\toprule
& What the agent does \\
\midrule
O1 & Never exploits \\
O2 & Reports the bug \\
O3 & Exploits in view, is corrected early, then complies \\
O4 & Deterred: exploits only where no correction can follow \\
\midrule
O0 & \textbf{Hides}: exploits everywhere except the audit window (the behaviour to remove) \\
\bottomrule
\end{tabularx}\par}

\textbf{Standing assumptions} (CLAUDE.md). The reward is updated from
trajectory pairs (PBRR), and PBRR's loss terms may be changed. An ideal reward
exists and every correction moves towards it. The overseer has perfect,
informative pairs.
\end{keybox}

% =====================================================================
\section{Why the agent hides}

For one round of $T$ steps, let $V_E$, $V_H$ and $V_0$ be the received return of
a round spent exploiting throughout, hiding (exploiting except in the audit
window), and acting honestly. The agent hides in round $k$ if and only if
hiding costs less than the correction it avoids:
\[
  V_E - V_H \;<\; p\,\Delta V_k,
  \qquad
  \Delta V_k = V_{\text{uncorrected}}(k{+}1{:}K) - V_{\text{corrected}}(k{+}1{:}K).
\]
That gives three levers:
\begin{itemize}
  \item \textbf{Motive}: shrink $\Delta V_k$, so being corrected costs nothing
    or is paid back. \hyperref[p1]{P1}, \hyperref[p2]{P2}, \hyperref[p3]{P3},
    \hyperref[p5]{P5}.
  \item \textbf{Means}: leave no predictable unwatched stretch to hide in.
    \hyperref[p4]{P4}.
  \item \textbf{Prize}: make exploiting lose reward in expectation, so there is
    nothing to hide. \hyperref[p6]{P6}.
\end{itemize}

\textbf{Evidence tags.}
\src{DP}: an exact solve of an idealised model (oracle patch plus a raise $c$
per step) on the Step 1 grid.
\src{Phase A}: the existing exact solve with the fitted neural $g$.
\src{closed form}: derived by hand.

% =====================================================================
\section{Summary}

\begin{center}
\small
\renewcommand{\arraystretch}{1.25}
\begin{tabularx}{\linewidth}{@{}l>{\raggedright\arraybackslash}p{2.7cm}lLLl@{}}
\toprule
& Proposal & Lever & What changes & Predicted optimum & End \\
\midrule
P1 & Drop regulariser terms & motive
   & PBRR loss without $L^-$ (or without any regulariser)
   & Corrected early, then complies; the size of the raise is uncontrolled \src{Phase A} & O3 \\
P2 & Set the size of the raise & motive
   & Regulariser target $0 \to c$
   & Hides below $c_1$; corrected early, then complies, above it \src{DP} & O3 \\
P3 & Don't credit the correction & motive
   & Counterfactual bootstrap, or myopic training plus approval (MONA)
   & Corrected early, then complies; the reward is unchanged \src{closed form} & O3 \\
P4 & Random hidden audits & means
   & Audited steps drawn at random, unseen by the agent
   & Never hides; a tiny raise makes early correction certain \src{DP} & O3/O4 \\
P5 & Bug bounty & motive
   & A \emph{report} action, verified by perfect pairs
   & Above $c_2$, reports at once and never exploits \src{DP} & O2 \\
P6 & Ex-post review & prize
   & Random whole-episode reviews; the correction is applied retroactively
   & Hides until $f \ge b/s$, then never exploits \src{DP} & O1 \\
\bottomrule
\end{tabularx}
\end{center}

% =====================================================================
\section{Shared evaluation}

\textbf{Pipeline.}
(1)~Phase A: the exact optimum with the real fitted $g$. No PPO until this
shows the intended optimum.
(2)~PPO on the grid, with the Step 2 config and 3 seeds, compared with step 1.
(3)~The best one or two proposals on the Step 3 suite.

\textbf{Metrics.}
\begin{itemize}
  \item \textbf{M1, end state}: which of O0--O4 the agent reaches. For O3,
    also the round of the first correction.
  \item \textbf{M2, hiding}: off-window exploitation with a clean audit window
    (Step 3's S1/S2). Target $\approx 0$.
  \item \textbf{M3, returns}: per-round received and true returns, on the
    Step 1 timeline.
  \item \textbf{M4, overshoot} (P1, P2, P5): received in a corrected round
    $\div$ $V_H$. Below 1, the agent hides. Above $V_E/V_H$, it seeks
    correction.
  \item \textbf{M5, motive probe}: compare the last round, where no correction
    can follow, with the earlier rounds. If they differ, the behaviour may be
    safe, but the motive is intact.
\end{itemize}

\textbf{Success} (pre-register). In $\ge 2$ of 3 seeds, the learned policy
matches its exact optimum on M1--M3, and that optimum is not O0.

% =====================================================================
\section{Proposals}

% ---------------------------------------------------------------------
\begin{proposal}{P1 \quad Drop regulariser terms}
\phantomsection\label{p1}\lever{motive}

\fld{Idea} The current correction is a fine. $L^-$ pins the preferred, honest
trajectory at $g=0$, so the fit can only make the exploit pay less. Without
that pin, the fit can also \emph{raise} honest behaviour. If a corrected future
pays more than a hidden one, the agent stops avoiding correction.

\fld{Change} Variants of Eq.~3:
(a) \code{no\_agree=skip};
(b) Eq.~1, with no regulariser at all;
(c) drop $L^-$ only;
(d) drop $L^+$ only.
(a) and (b) run today through \code{--set}. (c) and (d) need a flag in
\code{build\_design} and in \code{src/step3/pbrr.py}. If every rollout
exploits, the first correction has $|D^+|=0$, so (a)~=~(b)~=~(c) and
(d)~=~max1. The variants differ only from the second correction on.

\fld{Prediction} \src{Phase A, skip} The agent is corrected in round 0, then
complies. But nothing regulates the raise: on the grid, a corrected round pays
several times an exploit round. The raise sits on the reference behaviour, and
it vanishes at the second correction, when $L^+$ returns.

\fld{Experiments}
E1.1: Phase A for (a)--(d) and max1, over all corrections and 3 $g$ seeds.
Report M1, M3, M4 and where the raise lands.
E1.2: PPO for every variant whose optimum changes. Log $|D^+|$ at the first
correction: a single honest rollout brings $\lambda = 10/|D^+|$ back at its
largest, so skip may not fire in PPO.

\fld{Risks} The raise has no set size, so the agent may seek correction. It
lands where the reference goes, which in another environment can be a new
exploit. It does not survive later corrections.

\fld{Sources} \cite{pbrr, christiano}.
\end{proposal}

% ---------------------------------------------------------------------
\begin{proposal}{P2 \quad Set the size of the raise}
\phantomsection\label{p2}\lever{motive}

\fld{Idea} Keep the regulariser, but pull $g$ towards a constant $c$:
$\lambda \sum (g-c)^2$. The preference term cannot see a constant shift,
because all trajectories have length $T$. So the correction ranks trajectories
as before, and still moves towards the ideal reward; it just pays $c$ per step
from then on. This is the seamless-transition payment of~\cite{indifference},
written as a PBRR loss term.

\fld{Change} The regulariser's target in \code{build\_design} and in
\code{src/step3/pbrr.py}, with a new config field \code{g.reg\_target}.

\fld{Prediction} \src{DP} A corrected round pays $V_0 + cT$.
\begin{center}
\small
\begin{tabular}{@{}ll@{}}
\toprule
Raise & Optimum \\
\midrule
$c < c_1 = (V_H - V_0)/T$ & Hides, as now \\
$c_1 \le c < c_2 = (V_E - V_0)/T$ & Corrected early, then complies \\
$c \ge c_2$ & The same, but now prefers being corrected to never being audited \\
\bottomrule
\end{tabular}
\end{center}
$c_2$ is the exact-indifference payment $C = V^*(\text{old}) - V(\text{new})$.
The obvious objection, that the old reward is still higher so the agent should
lean towards it, is answered by $c$: it pays back exactly that gap.

What remains is a precision problem. Any underpayment restores hiding in full.
The safe band $[c_1, c_2)$ is $(V_E - V_H)/T$ wide, the per-step cost of hiding,
so it is narrow when hiding is cheap. A wider audit window makes hiding
costlier and lowers $c_1$; P4 is the limit of this.

\fld{Experiments}
E2.1: Phase A with the fitted $g$, at $c/c_1 \in \{0, 0.9, 0.99, 1.01, 1.1\}$ and
at $c_2$. Check that the neural $g$ adds $c$ everywhere, since spillover can
move the threshold.
E2.2: phase diagrams over $c \times W$ and $c \times p$.
E2.3: PPO at $c \in \{0,\, 0.9c_1,\, 1.1c_1,\, 1.1c_2\}$.
E2.4: Step 3, with $c_1$ computed for each environment from its exact solver.

\fld{Risks} Setting $c$ needs $V_H$, i.e.\ the bug's value, for each
environment. Underpaying brings back full hiding; overpaying produces
correction-seeking.

\fld{Sources} \cite{indifference} (Sec.~5, Def.~16, Thm.~17), \cite{pbrr}.
\end{proposal}

% ---------------------------------------------------------------------
\begin{proposal}{P3 \quad Don't credit the correction: counterfactual critic, or MONA}
\phantomsection\label{p3}\lever{motive}

\fld{Idea} Leave the reward alone and change how the learner values a
correction. If the learner gives no credit to what the correction does to its
future, then being corrected costs it nothing and there is nothing to protect.

\fld{Change}
(a) \emph{Counterfactual critic.} At a correcting boundary, bootstrap with
$V(o'\text{ with the old correction count})$. Train that value on a few shadow
envs (e.g.\ 4 of 16) where corrections are switched off.
(b) \emph{MONA.} Train myopically, with each round terminal for the learner.
Give legitimate long-term value back as the overseer's approval of the round's
end state. The Step 1 harness has \code{myopic}; Steps 2 and 3 need it added.

\fld{Prediction} \src{closed form} At every correction level, the agent plays
that level's no-overseer optimum: it exploits openly, is corrected, and then
complies. Step 1 already showed that the myopic control removes hiding, both in
DP and in PPO.

\fld{Control} On the grid, (a), (b) and plain myopia coincide. Add a
\emph{bank} tile: visits in round $k$ pay in round $k{+}1$, under both rewards.
Then (a) uses the bank, plain myopia does not, and (b) uses it only if the
approval credits it.

\fld{Experiments}
E3.1: DP at each correction level, with and without the bank.
E3.2: PPO, comparing (a), (b), plain myopia and the baseline. Run (b) first,
since it is cheaper.
E3.3: the Step 3 suite.

\fld{Risks} A critic error is a payment error. For (b), the approval must
judge long-term value well. Indifference is not corrigibility: the agent won't
help find the bug. Both variants change the learner, which assumes we control
its training objective.

\fld{Sources} \cite{tampering} (current-RF, TI-ignoring), \cite{pathspecific},
\cite{indifference} (Sec.~3), \cite{mona}.
\end{proposal}

% ---------------------------------------------------------------------
\begin{proposal}{P4 \quad Random hidden audits}
\phantomsection\label{p4}\lever{means}

\fld{Idea} Hiding needs a predictable unwatched stretch. Audit random steps and
never tell the agent which ones. Then no part of the round is safe.

\fld{Change} Each step is audited with probability $q$; $q = W/T$ keeps today's
budget. The agent is caught at the boundary if any audited step collected the
bonus.
\emph{Variant (b).} Replay the current policy in a sandbox on whole rounds. If
the sandbox can't be told apart from training, this is $W = T$, and the trigger
no longer depends on when the agent exploits (unriggable, as in
\cite{pitfalls}).

\fld{Prediction} \src{DP} The agent never hides, at any $q$.
\begin{itemize}
  \item \textbf{Low $q$}: it exploits openly until caught. A forward-only
    correction costs at most about one exploit round in expectation, and the
    agent earns that round in the round it is caught.
  \item \textbf{High $q$}: deterrence (O4) and early correction (O3) tie,
    because each keeps one exploit round. Any raise $c > 0$ picks O3.
\end{itemize}
With random audits, then, P2's $c_1$ falls to about 0, and the bug's value
need not be known. Step 1's no-clock condition is the information-only version
of this idea.

\fld{Experiments}
E4.1: Phase A at $q \in \{0.1, 0.5, 1, 2\} \times W/T$, with
$c \in \{0, \text{small}\}$.
E4.2: PPO at $q = W/T$, with $c \in \{0, \text{small}\}$.
E4.3: motive test. Add a second exploit with a smaller bonus that is audited at
$q/10$. With $c = 0$, the agent should move to it if the motive survives.
E4.4: Step 3.

\fld{Risks} It removes the means, not the motive (E4.3). At low $q$, the agent
exploits for longer before it is corrected. Per-step audits can miss
multi-step pattern exploits (Boat Race).

\fld{Sources} \cite{inspection, decoupled, pitfalls}.
\end{proposal}

% ---------------------------------------------------------------------
\begin{proposal}{P5 \quad Bug bounty: let the agent report the bug}
\phantomsection\label{p5}\lever{motive}

\fld{Idea} Add a \emph{report} action. If perfect pairs confirm the bug, the
overseer corrects at once and pays $c$ per step from then on. The agent can
then get corrected without exploiting in view, or without exploiting at all.

\fld{Change} A report action, pressed on the suspected state. On a report, the
overseer builds pairs (the reported behaviour vs the reference), fits PBRR, and
applies the result \emph{immediately}, which is a protocol change. Only
verified reports are paid.

\fld{Prediction} \src{DP, immediate correction}
For $c < c_1$, the agent never reports and hides.
For $c_1 \le c < c_2$, it exploits the first round and is then corrected.
For $c \ge c_2$, it reports at the start and never exploits.
\src{closed form} The payment must come per step after the correction. With a
lump sum, the agent exploits first and reports at the end of the round.

\fld{Experiments}
E5.1: DP with the fitted $g$, at $c$ around $c_1$ and $c_2$, with
$p \in \{0.5, 1\}$.
E5.2: PPO.
E5.3: a bug that changes every episode, so the agent must find it before
reporting it.
E5.4: a fee for unverified reports, to check that the agent doesn't spam them.

\fld{Risks} It needs $c \ge c_2$, which means correction-seeking by design. That
is safe only because verification is perfect. With a fixed bug the task is
trivial (E5.3).

\fld{Sources} \cite{confessions, corrigibility}.
\end{proposal}

% ---------------------------------------------------------------------
\begin{proposal}{P6 \quad Ex-post review: apply the correction retroactively}
\phantomsection\label{p6}\lever{prize}

\fld{Idea} A forward-only correction can take away only future exploitation,
which is why P4's audits do not deter. So let the review reach backwards.
After each episode, a random fraction $f$ of episodes is reviewed in full, and
the PBRR correction is applied to every step of each reviewed episode. Hiding
then protects nothing, and if exploiting loses reward in expectation, there is
nothing to hide.

\fld{Change} At episode end, with probability $f$, the overseer builds pairs
from the whole episode, fits PBRR, and adds (corrected $-$ received) as a
terminal reward. The agent never knows which episodes are reviewed. The window
audit and the forward corrections stay as they are.

\fld{Prediction} \src{DP} In expectation, the bonus shrinks to
$b_{\text{eff}} = b - f s$, where $s$ is how much the retroactive correction
takes off one exploit step. While $b_{\text{eff}} > 0$, the agent hides exactly
as before, only for smaller stakes. Once $b_{\text{eff}} \approx 0$, it stops
exploiting altogether. The threshold is therefore $f^* = b/s$. A clawback of the
bonus alone ($s = b$) needs $f^* = 1$. PBRR's fine overshoots the bonus
($s > b$), which pushes $f^*$ well below 1.

\fld{Experiments}
E6.1: Phase A at $f/f^* \in \{0.5, 0.9, 1.1, 1.5\}$.
E6.2: PPO. The penalty arrives up to an episode after the exploit, so check
that PPO carries the credit that far.
E6.3: Step 3, where $s$ differs per environment.

\fld{Risks} This is deterrence: the motive is intact (M5). The agent resumes if
$f$ drops, or if it can tell which episodes are reviewed. The overseer learns
nothing new. $f^*$ rests on the size of PBRR's fine, which nothing controls.

\fld{Sources} \cite{becker, inspection}.
\end{proposal}

% =====================================================================
\section{Considered and dropped}
\begin{itemize}
  \item \textbf{Penalising different behaviour while monitored.} The clock
    carries the same information as the monitor bit, so the agent would
    re-route through the clock.
  \item \textbf{Forward-only deterrence.} It cannot deter at a low audit rate
    (P4). Deterrence needs a retroactive fine (P6).
\end{itemize}

% =====================================================================
\section{Suggested order and open questions}
\begin{enumerate}
  \item \textbf{One Phase A sweep, no PPO}: P1's variants and the sweeps over
    $c$ (P2), $q \times c$ (P4) and $f$ (P6). Port the idealised DP into
    \code{CheckLog}.
  \item \textbf{PPO on the grid}: P2, P4, P6, and P1 as Eq.~1.
  \item \textbf{P3, then P5}: both need new code (the myopic condition, the
    critic bootstrap and shadow envs, the report action) and the bank-tile
    variant.
  \item \textbf{Step 3 suite} for the best one or two proposals.
\end{enumerate}

\textbf{Open questions.}
\begin{itemize}
  \item How should the end states rank? This draft uses
    O1 $\approx$ O2 $>$ O3 $>$ O4 $>$ O0.
  \item P1: drop only $L^-$, or the whole regulariser?
  \item Is immediate correction on a verified report (P5) an acceptable
    protocol change?
  \item P3 changes the learner. Do we control the training objective, or only
    the overseer?
  \item Can the overseer replay the policy in a sandbox (P4b), or review whole
    episodes after the fact (P6)?
\end{itemize}

% =====================================================================
\begin{thebibliography}{99}
\small
\setlength{\itemsep}{2pt}

\bibitem{indifference}
S.~Armstrong and X.~O'Rourke.
\newblock `Indifference' methods for managing agent rewards.
\newblock arXiv:1712.06365, 2017.

\bibitem{pitfalls}
S.~Armstrong, J.~Leike, L.~Orseau, and S.~Legg.
\newblock Pitfalls of learning a reward function online.
\newblock In \emph{IJCAI}, 2020. arXiv:2004.13654.

\bibitem{inspection}
R.~Avenhaus, B.~von Stengel, and S.~Zamir.
\newblock Inspection games.
\newblock In \emph{Handbook of Game Theory with Economic Applications}, vol.~3,
  ch.~51. Elsevier, 2002.

\bibitem{becker}
G.~S.~Becker.
\newblock Crime and punishment: an economic approach.
\newblock \emph{Journal of Political Economy}, 76(2):169--217, 1968.

\bibitem{christiano}
P.~Christiano, J.~Leike, T.~Brown, M.~Martic, S.~Legg, and D.~Amodei.
\newblock Deep reinforcement learning from human preferences.
\newblock In \emph{NeurIPS}, 2017. arXiv:1706.03741.

\bibitem{tampering}
T.~Everitt, M.~Hutter, R.~Kumar, and V.~Krakovna.
\newblock Reward tampering problems and solutions in reinforcement learning: a
  causal influence diagram perspective.
\newblock \emph{Synthese}, 2021. arXiv:1908.04734.

\bibitem{mona}
S.~Farquhar, V.~Varma, D.~Lindner, D.~Elson, C.~Biddulph, I.~Goodfellow, and
  R.~Shah.
\newblock {MONA}: myopic optimization with non-myopic approval can mitigate
  multi-step reward hacking.
\newblock arXiv:2501.13011, 2025.

\bibitem{pathspecific}
S.~Farquhar, R.~Carey, and T.~Everitt.
\newblock Path-specific objectives for safer agent incentives.
\newblock In \emph{AAAI}, 2022. arXiv:2204.10018.

\bibitem{pbrr}
S.~Hatgis-Kessell, L.~M.~Bhamidipaty, and E.~Brunskill.
\newblock Repairing reward functions with feedback to mitigate reward hacking.
\newblock arXiv:2510.13036v2, 2025.

\bibitem{confessions}
M.~Joglekar, J.~Chen, G.~Wu, J.~Yosinski, J.~Wang, B.~Barak, and A.~Glaese.
\newblock Training {LLMs} for honesty via confessions.
\newblock arXiv:2512.08093, 2025.

\bibitem{corrigibility}
N.~Soares, B.~Fallenstein, E.~Yudkowsky, and S.~Armstrong.
\newblock Corrigibility.
\newblock In \emph{AAAI Workshop on AI and Ethics}, 2015.

\bibitem{decoupled}
J.~Uesato, R.~Kumar, V.~Krakovna, T.~Everitt, R.~Ngo, and S.~Legg.
\newblock Avoiding tampering incentives in deep {RL} via decoupled approval.
\newblock arXiv:2011.08827, 2020.

\end{thebibliography}

\end{document}
